% ============================================================
%  chapters/ch01-introduction.tex
% ============================================================

\chapter{Introduction}
\label{ch:intro}

\section{Background and Motivation}
\label{sec:bg}

% Replace with your own introduction text.
Medical imaging is essential for modern diagnosis, treatment planning, and
disease monitoring.  X-ray and computed tomography (CT) are the most widely
used modalities because they are fast, accessible, and cost-effective.  In
low-income countries with limited access to MRI or ultrasound, X-ray remains
the primary diagnostic tool for conditions ranging from fractures to pneumonia.

However, image quality is often degraded during acquisition, reducing
diagnostic accuracy.  This can lead to missed findings, false positives, or
unnecessary repeat scans.  Poor image quality also forces clinicians to request
additional views or higher radiation doses, increasing patient risk.

\section{Problem Statement}
\label{sec:ps}

% Summarise the core research problem.

\section{Research Objectives}
\label{sec:obj}

% List your research objectives.

\section{Scope and Limitations}
\label{sec:scope}

% Define the scope of your work.

\section{Organisation of the Thesis}
\label{sec:org}

The remainder of this thesis is organised as follows.
Chapter~\ref{ch:litreview} presents a comprehensive literature review.
Chapter~\ref{ch:methodology} describes the proposed methodology.
Chapter~\ref{ch:results} presents experimental results and discussion.
Chapter~\ref{ch:conclusion} concludes the thesis and outlines future work.
